% section: 数学的帰納法

  \section{数学的帰納法}\label{節：数学的帰納法}
  \noindent\textbf{【本編で使う場面】}\par
  数項からの予想をすべての項の証明へ進める方法を学ぶ.\sectionref*{節：漸化式を解く前に}の検算や,\sectionref*{節：虚数解型隣接3項間漸化式}の\bookterm{general}{一般項}の確認で使う.

    ここまでは計算に用いる公式や式変形を扱った. ここでは,自然数に関する主張を証明する方法を説明する.
    自然数$n$についての主張$P(n)$が,すべての$n\ge1$で成り立つことを示したいとする.
    このとき,以下の2つを示せばよい.
    \begin{enumerate}[label=(\roman*),parsep=3pt]
      \item $P(1)$が成り立つ.
      \item $n\ge1$を任意にとり,\,$P(n)$が成り立つと仮定すると,\,$P(n+1)$も成り立つ.
    \end{enumerate}
    この2つが示されれば,\,(i)より$P(1)$が成り立ち,\,(ii)より$P(1)\Rightarrow P(2)$,\,$P(2)\Rightarrow P(3)$,\,$\ldots$と次々に成り立つことが分かるので,すべての$n\ge1$について$P(n)$が成り立つ.
    ドミノ倒しに例えられることが多い.\,(i)は最初のドミノを倒すこと,\,(ii)は「どのドミノについても,それが倒れれば次のドミノも倒れる」ように並べることに対応する.

    \sectionref*{節：極形式と複素平面}では,この方法を用いて\bookterm{de-moivre}{ド・モアブルの定理}を証明する.
    また,\,\chapterref{章：漸化式の解説}では,\bookterm{recurrence}{漸化式}から予想した\bookterm{general}{一般項}の形を確かめる場面や,\bookterm{trigonometric}{三角関数}の\bookterm{addition}{加法定理}から\bookterm{recurrence}{漸化式}を導く場面(\sectionref{節：漸化式を解く前に}や\sectionref{節：虚数解型隣接3項間漸化式}など)で\bookterm{induction}{数学的帰納法}を用いる.
